\Scene{2}[The Forest of Arden]

\enter{\textsc{Orlando,} with his arm in a sling, and \textsc{Oliver}}


\begin{prose_speech}[Orlando] 
 Is 't possible that on so little acquaintance you should like her? That, but seeing, you should love her? And loving, woo? And wooing, she should grant? And will you persever to enjoy her? 
 \end{prose_speech}

\begin{prose_speech}[Oliver] 
 Neither call the giddiness of it in question, the poverty of her, the small acquaintance, my sudden wooing, nor her sudden consenting, but say with me <I love Aliena>; say with her that she loves me; consent with both that we may enjoy each other. It shall be to your good, for my father's house and all the revenue that was old Sir Rowland's will I estate upon you, and here live and die a shepherd. 
 \end{prose_speech}

\enter{\textsc{Rosalind,} as \textsc{Ganymede}}


\begin{prose_speech}[Orlando] 
 You have my consent. Let your wedding be tomorrow. Thither will I invite the Duke and all 's contented followers. Go you and prepare Aliena, for, look you, here comes my Rosalind. 
 \end{prose_speech}

\verseline[Rosalind as Ganymede]{\stage{To \textsc{Oliver}} God save you, brother.}
\begin{prose_speech}[Oliver] 
 And you, fair sister.	
 \end{prose_speech}
 \exit{\textsc{Oliver}}

\begin{prose_speech}[Rosalind as Ganymede] 
   O my dear Orlando, how it grieves me to see thee wear thy heart in a scarf. 
 \end{prose_speech}

\begin{prose_speech}[Orlando] 
 It is my arm. 
 \end{prose_speech}

\begin{prose_speech}[Rosalind as Ganymede] 
   I thought thy heart had been wounded with the claws of a lion. 
 \end{prose_speech}

\begin{prose_speech}[Orlando] 
 Wounded it is, but with the eyes of a lady. 
 \end{prose_speech}

\begin{prose_speech}[Rosalind as Ganymede] 
   Did your brother tell you how I counterfeited to swoon when he showed me your handkercher? 
 \end{prose_speech}

\begin{prose_speech}[Orlando] 
 Ay, and greater wonders than that. 
 \end{prose_speech}

\begin{prose_speech}[Rosalind as Ganymede] 
   O, I know where you are. Nay, 'tis true. There was never anything so sudden but the fight of two rams, and Caesar's thrasonical brag of <I came, saw, and overcame.> For your brother and my sister no sooner met but they looked, no sooner looked but they loved, no sooner loved but they sighed, no sooner sighed but they asked one another the reason, no sooner knew the reason but they sought the remedy; and in these degrees have they made a pair of stairs to marriage, which they will climb incontinent, or else be incontinent before marriage. They are in the very wrath of love, and they will together. Clubs cannot part them. 
 \end{prose_speech}

\begin{prose_speech}[Orlando] 
 They shall be married tomorrow, and I will bid the Duke to the nuptial. But O, how bitter a thing it is to look into happiness through another man's eyes. By so much the more shall I tomorrow be at the height of heart-heaviness by how much I shall think my brother happy in having what he wishes for. 
 \end{prose_speech}

\begin{prose_speech}[Rosalind as Ganymede] 
   Why, then, tomorrow I cannot serve your turn for Rosalind? 
 \end{prose_speech}

\begin{prose_speech}[Orlando] 
 I can live no longer by thinking. 
 \end{prose_speech}

\begin{prose_speech}[Rosalind as Ganymede] 
   I will weary you then no longer with idle talking. Know of me then—for now I speak to some purpose—that I know you are a gentleman of good conceit. I speak not this that you should bear a good opinion of my knowledge, insomuch I say I know you are. Neither do I labour for a greater esteem than may in some little measure draw a belief from you to do yourself good, and not to grace me. Believe then, if you please, that I can do strange things. I have, since I was three year old, conversed with a magician, most profound in his art and yet not damnable. If you do love Rosalind so near the heart as your gesture cries it out, when your brother marries Aliena shall you marry her. I know into what straits of fortune she is driven, and it is not impossible to me, if it appear not inconvenient to you, to set her before your eyes tomorrow, human as she is, and without any danger. 
 \end{prose_speech}

\begin{prose_speech}[Orlando] 
 Speak'st thou in sober meanings? 
 \end{prose_speech}

\begin{prose_speech}[Rosalind as Ganymede] 
   By my life I do, which I tender dearly, though I say I am a magician. Therefore put you in your best array, bid your friends; for if you will be married tomorrow, you shall, and to Rosalind, if you will. 

\enter{\textsc{Silvius and Phoebe}}

Look, here comes a lover of mine and a lover of hers.
 \end{prose_speech}

\begin{verse_speech}[Phoebe]
\stage{To \textsc{Rosalind}}
Youth, you have done me much ungentleness\\
To show the letter that I writ to you.
\end{verse_speech}

\begin{verse_speech}[Rosalind as Ganymede]
I care not if I have. It is my study\\
To seem despiteful and ungentle to you.\\
You are there followed by a faithful shepherd.\\
Look upon him, love him; he worships you.
\end{verse_speech}

\begin{verse_speech}[Phoebe] 
\stage{To \textsc{Silvius}}
Good shepherd, tell this youth what 'tis to love.
\end{verse_speech}

\begin{verse_speech}[Silvius]
	It is to be all made of sighs and tears,\\
And so am I for Phoebe.
\end{verse_speech}

\verseline[Phoebe]{And I for Ganymede.}
\verseline[Orlando]{And I for Rosalind.}
\verseline[Rosalind as Ganymede]{  And I for no woman.}
\begin{verse_speech}[Silvius]
	It is to be all made of faith and service,\\
And so am I for Phoebe.
\end{verse_speech}

\verseline[Phoebe]{And I for Ganymede.}
\verseline[Orlando]{And I for Rosalind.}
\verseline[Rosalind as Ganymede]{  And I for no woman.}
\begin{verse_speech}[Silvius]
	It is to be all made of fantasy,\\
All made of passion and all made of wishes,\\
All adoration, duty, and observance,\\
All humbleness, all patience and impatience,\\
All purity, all trial, all observance,\\
And so am I for Phoebe.
\end{verse_speech}

\verseline[Phoebe]{And so am I for Ganymede.}
\verseline[Orlando]{And so am I for Rosalind.}
\verseline[Rosalind as Ganymede]{  And so am I for no woman.}
\verseline[Phoebe]{If this be so, why blame you me to love you?}
\verseline[Silvius]{If this be so, why blame you me to love you?}
\verseline[Orlando]{If this be so, why blame you me to love you?}
\verseline[Rosalind as Ganymede]{Why do you speak too, <Why blame you me to love you?>}
\verseline[Orlando]{To her that is not here, nor doth not hear.}
\begin{prose_speech}[Rosalind as Ganymede] 
   Pray you, no more of this. 'Tis like the howling of Irish wolves against the moon. \stage{To \textsc{Silvius}} I will help you if I can. 
   \stage{To \textsc{Phoebe}} I would love you if I could.
   
   —Tomorrow meet me all together. 
   \stage{To \textsc{Phoebe}} I will marry you if ever I marry woman, and I'll be married tomorrow. 
   \stage{To \textsc{Orlando}} I will satisfy you if ever I satisfy man, and you shall be married tomorrow. 
   \stage{To \textsc{Silvius}} I will content you, if what pleases you contents you, and you shall be married tomorrow. 
   \stage{To \textsc{Orlando}} As you love Rosalind, meet. 
   \stage{To \textsc{Silvius}} As you love Phoebe, meet.—And as I love no woman, I'll meet. So fare you well. I have left you commands. 
 \end{prose_speech}

\verseline[Silvius]{I'll not fail, if I live.}
\verseline[Phoebe]{Nor I.}
\verseline[Orlando]{Nor I.}\exeunt{}